\booksection{Chapitre 11 Annexe : Connaissances en technologies d'image professionnelles}

\chapterintro{Pour vraiment comprendre les secrets de la photographie et de la vidéo sur smartphone, vous devez connaître quelques notions professionnelles de technologie d'image. Ces connaissances vous aident non seulement à mieux utiliser l'appareil photo du smartphone, mais aussi à être plus critique dans le choix d'un équipement.}

\booksubsection{Résolution, fréquence d'images et débit : les trois piliers de la qualité vidéo}

\subsubsection*{Résolution :}

La résolution détermine la netteté de l'image, exprimée en nombre de pixels.

\begin{itemize}
  \item \textbf{720p (HD)}: 1280×720 pixels, netteté d'entrée de gamme.
  \item \textbf{1080p (Full HD)}: 1920×1080 pixels, la résolution la plus courante actuellement.
  \item \textbf{2K (QHD)}: 2560×1440 pixels, plus net que le 1080p.
  \item \textbf{4K (UHD)}: 3840×2160 pixels, ultra haute définition, riche en détails.
  \item \textbf{8K}: 7680×4320 pixels, norme future, prise en charge limitée actuellement.
\end{itemize}

Plus la résolution est élevée, plus l'image est nette, mais le fichier est plus volumineux et les exigences de stockage et de traitement sont plus élevées.

\subsubsection*{Fréquence d'images (Frame Rate) :}

La fréquence d'images détermine la fluidité de l'image, exprimée en images par seconde (fps).

\begin{itemize}
  \item \textbf{24 fps}: Fréquence standard du cinéma, avec un « effet cinématographique » unique.
  \item \textbf{30 fps}: Courant pour la télévision et les vidéos web, fluide et naturel.
  \item \textbf{60 fps}: Haute fréquence, adapté aux scènes sportives et à l'enregistrement de jeux, très fluide.
  \item \textbf{120 fps/240 fps}: Très haute fréquence, utilisable pour les vidéos au ralenti.
\end{itemize}

Plus la fréquence est élevée, plus l'image est fluide, meilleur est l'effet au ralenti, mais le fichier est plus volumineux.

\subsubsection*{Débit (Bitrate) :}

Le débit détermine la densité des données de la vidéo, exprimé en volume de données par seconde (Mbps).

\begin{itemize}
  \item Plus le débit est élevé, meilleure est la qualité d'image, plus les détails sont préservés, mais le fichier est plus volumineux.
  \item Un débit trop bas provoque l'apparition de mosaïques, de blocs de couleur et autres artefacts de compression.
  \item Le tournage vidéo professionnel nécessite généralement un débit élevé pour préserver davantage de marge de manœuvre en post-production.
\end{itemize}

\booksubsection{Formats de codage vidéo : H.264, H.265 et ProRes}

\subsubsection*{H.264 (AVC) :}

\begin{itemize}
  \item \textbf{Caractéristiques}: Le format de codage le plus largement utilisé, excellente compatibilité.
  \item \textbf{Avantage}: Pratiquement tous les appareils peuvent le lire, efficacité de compression moyenne.
  \item \textbf{Inconvénient}: Comparé aux nouveaux formats, efficacité de compression plus faible, fichiers plus volumineux.
  \item \textbf{Scénario d'utilisation}: Partage de vidéos en ligne, tournage quotidien pour le grand public.
\end{itemize}

\subsubsection*{H.265 (HEVC) :}

\begin{itemize}
  \item \textbf{Caractéristiques}: Version améliorée du H.264, efficacité de compression améliorée d'environ 50 \%.
  \item \textbf{Avantage}: Fichier plus petit à qualité égale, prend en charge des résolutions et fréquences d'images plus élevées.
  \item \textbf{Inconvénient}: Le codage et le décodage nécessitent un processeur plus puissant, certains appareils anciens ne sont pas compatibles.
  \item \textbf{Scénario d'utilisation}: Tournage vidéo 4K, scénarios avec espace de stockage limité.
\end{itemize}

\subsubsection*{Apple ProRes :}

\begin{itemize}
  \item \textbf{Caractéristiques}: Codec vidéo professionnel développé par Apple, compression quasi sans perte.
  \item \textbf{Avantage}: Qualité d'image extrêmement élevée, grande marge de post-production, couleurs précises.
  \item \textbf{Inconvénient}: Fichiers extrêmement volumineux, nécessite beaucoup d'espace de stockage, le traitement nécessite un équipement professionnel.
  \item \textbf{Scénario d'utilisation}: Production vidéo professionnelle, tournage de niveau cinématographique.
\end{itemize}

\subsubsection*{ProRes RAW :}

\begin{itemize}
  \item \textbf{Caractéristiques}: Combine les avantages du ProRes et du format RAW.
  \item \textbf{Qu'est-ce que le RAW}: Données brutes du capteur non traitées, préservant la quantité maximale d'informations.
  \item \textbf{Avantage}: Marge de réglage en post-production extrêmement grande, possibilité de modifier l'exposition, la balance des blancs et d'autres paramètres sans perte de qualité d'image.
  \item \textbf{Inconvénient}: Fichiers extrêmement volumineux, nécessite un logiciel de post-production professionnel.
  \item \textbf{Scénario d'utilisation}: Production professionnelle haut de gamme. À partir de la série iPhone 15 Pro, la vidéo Apple ProRes est prise en charge (jusqu'à 4K 60 fps) ; l'iPhone 17 Pro prend en charge en plus l'enregistrement ProRes RAW directement dans l'appareil.
\end{itemize}

\booksubsection{HDR : technologie à haute plage dynamique}

La technologie HDR (High Dynamic Range) permet de préserver simultanément les détails dans les parties les plus claires et les plus sombres de l'image.

\begin{itemize}
  \item \textbf{HDR10}: La norme HDR la plus basique, bonne compatibilité.
  \item \textbf{HDR10+}: standard promu par Samsung, prend en charge les métadonnées dynamiques.
  \item \textbf{Dolby Vision}: le meilleur standard HDR, développé par Dolby, offre la meilleure qualité d'image, mais nécessite une licence.
  \item \textbf{HLG}: standard HDR de qualité broadcast, compatible avec les écrans ordinaires.
\end{itemize}

Lors du tournage de vidéos HDR, vous avez besoin d'un écran prenant en charge le HDR pour voir l'effet complet.

\booksubsection{IV. Espace colorimétrique et profondeur de couleur}

\subsubsection*{Espace colorimétrique :}

\begin{itemize}
  \item \textbf{sRGB}: espace colorimétrique standard, la meilleure compatibilité.
  \item \textbf{Display P3}: standard à large gamut promu par Apple, des couleurs plus riches.
  \item \textbf{Rec.2020}: espace colorimétrique standard pour les vidéos HDR, le gamut le plus large.
\end{itemize}

\subsubsection*{Profondeur de couleur :}

\begin{itemize}
  \item \textbf{8-bit}: 16,7 millions de couleurs, standard vidéo ordinaire.
  \item \textbf{10-bit}: 1,07 milliard de couleurs, standard vidéo professionnel, transitions de couleurs plus fluides.
  \item \textbf{12-bit}: 68 milliards de couleurs, standard professionnel haut de gamme.
\end{itemize}

\begin{tcbinfo}{Conseils pratiques}
\begin{itemize}
  \item Usage quotidien : 1080P/4K + 30fps + encodage H.265, équilibre entre qualité d'image et taille de fichier.
  \item Filmage de sport : 1080P/4K + 60fps, garantit la fluidité.
  \item Création professionnelle : 4K + ProRes (si pris en charge), préserve un maximum de marge de post-production.
  \item Conseil de stockage : les vidéos à haut débit occupent beaucoup d'espace, il est recommandé de les sauvegarder régulièrement sur ordinateur ou cloud.
\end{itemize}

\end{tcbinfo}

\begin{tcbpractice}{À explorer : vos paramètres vidéo sur smartphone}
\begin{enumerate}
  \item Ouvrez l'application appareil photo, passez en mode vidéo.
  \item Consultez les options de résolution et de fréquence d'images, essayez différentes combinaisons.
  \item Si pris en charge, vérifiez s'il existe des options avancées comme HDR, ProRes, etc.
  \item Filmez quelques séquences avec différents réglages, comparez la qualité d'image et la taille des fichiers.
\end{enumerate}

\end{tcbpractice}

\booksubsection{V. Comparaison entre RAW, HEIF et JPEG}

\subsubsection*{Qu'est-ce que le format RAW :}

Le format RAW contient les « données brutes » enregistrées par le capteur de l'appareil photo — des informations non compressées et non traitées. C'est comme des ingrédients crus qui nécessitent une « cuisson » en post-traitement pour devenir un plat fini. Les smartphones, en raison de leurs limites de stockage et de puissance de calcul, produisent par défaut des JPEG ou HEIF déjà « cuits » ; le RAW doit être activé spécifiquement.

\subsubsection*{Comparaison RAW vs HEIF vs JPEG :}

\begin{itemize}
  \item \textbf{JPEG}: Le format le plus ancien, meilleure compatibilité (ouvrable sur tous les appareils). Profondeur de couleur 8 bits, compression avec perte, fichiers les plus petits (2-5 Mo par photo). Espace de post-traitement limité ; éclaircir fortement les ombres fait apparaître des artefacts.
  \item \textbf{HEIF/HEIC}: Nouveau format promu par Apple ; à qualité égale, la taille du fichier est la moitié de celle du JPEG (1-3 Mo par photo). Prend en charge la profondeur de couleur 10 bits et le canal alpha. Inconvénient : certains anciens appareils et Windows nécessitent un plugin pour l'ouvrir.
  \item \textbf{RAW}: Données brutes non compressées, fichiers les plus volumineux (20-100 Mo par photo). Profondeur de couleur généralement 12-16 bits, espace de post-traitement maximal. Inconvénient : nécessite un logiciel professionnel (Lightroom, Snapseed, etc.) pour ouvrir et traiter, ne peut pas être partagé directement.
\end{itemize}

\subsubsection*{Pourquoi les photographes professionnels utilisent le RAW :}

\begin{itemize}
  \item \textbf{Balance des blancs ajustable sans perte en post-traitement}: Si la balance des blancs est mauvaise en JPEG (par exemple, jaunâtre sous une ampoule à incandescence), la correction en post-traitement dégrade la qualité ; le RAW permet de redéfinir la température de couleur sans aucune perte.
  \item \textbf{Grande tolérance d'exposition}: En JPEG, les ombres sous-exposées de 2 IL sont complètement noires, et l'éclaircissement ne révèle que du bruit ; le RAW peut « sauver » des détails utilisables dans des ombres sous-exposées de 3-4 IL. Même principe pour la surexposition.
  \item \textbf{Plus de détails dans les ombres et les hautes lumières}: La plage dynamique du RAW est bien supérieure à celle du JPEG, offant beaucoup plus de marge pour des hautes lumières non brûlées et des ombres non bouchées.
  \item \textbf{Netteté, réduction du bruit et correction d'objectif contrôlables en post-traitement}: Le traitement intégré à l'appareil est « définitif » ; le RAW permet des ajustements répétés, toujours à partir des données d'origine.
\end{itemize}

\subsubsection*{Comment activer le RAW sur un smartphone :}

\begin{itemize}
  \item \textbf{iPhone ProRAW}: Uniquement sur iPhone 12 Pro et modèles Pro supérieurs. Activez « Apple ProRAW » dans « Réglages Appareil photo Formats ». Un bouton « RAW » apparaît alors en haut de l'interface de l'appareil photo ; touchez-le pour l'activer. Environ 25 Mo par photo, à activer seulement en cas de besoin.
  \item \textbf{Android Motion Raw / Ultra RAW}: Le nom varie selon les marques. Sur les modèles phares Xiaomi/OPPO/vivo, l'option RAW se trouve dans le « Mode pro » ou « Plus » de l'appareil photo ; Huawei l'active dans le « Mode pro » ; Samsung Galaxy S series propose une option RAW dans le mode « Pro ». Taille de fichier généralement 20-50 Mo.
  \item \textbf{Conseil général}: Pas besoin d'activer le RAW pour les photos quotidiennes, les fichiers sont trop volumineux ; activez-le uniquement pour les scènes importantes (paysages, portraits, post-traitement nécessaire).
\end{itemize}

\booksubsection{VI. Principes de la photographie computationnelle}

La photographie moderne sur smartphone ne peut plus se passer de la « photographie computationnelle » — c'est-à-dire s'appuyer sur des algorithmes et pas seulement sur l'objectif et le capteur pour produire des photos. Comprendre ces principes fondamentaux vous aidera à mieux saisir ce que fait le smartphone au moment où vous appuyez sur le déclencheur.

\subsubsection*{Fusion multi-images :}

Lorsque vous appuyez une fois sur le déclencheur, le smartphone prend en réalité plusieurs photos. L'algorithme aligne et fusionne pondérément les images pour obtenir une photo finale de haute qualité.

\begin{itemize}
  \item \textbf{Mode nuit}: Prend plusieurs photos sous-exposées en rafale pendant 1-3 secondes, les aligne puis les superpose — le bruit s'annule et la luminosité s'accumule. Résultat : des images nettes même en faible lumière.
  \item \textbf{HDR (High Dynamic Range)}: Prend trois (ou plusieurs) photos sous-exposée, normalement exposée et surexposée, puis les fusionne en une seule photo avec des détails dans les hautes lumières et les ombres. Particulièrement efficace en contre-jour.
  \item \textbf{Deep Fusion (Apple) / Photonic Engine}: Pour les scènes de faible à moyenne luminosité, analyse et fusionne plusieurs images pixel par pixel, en améliorant particulièrement les textures et les détails. Les détails comme les tricots et les cheveux seront nettement plus clairs.
\end{itemize}

\subsubsection*{Segmentation sémantique :}

L'IA ne traite pas l'image entière de manière uniforme ; elle « comprend » d'abord ce qu'il y a dans l'image — ciel, visages, herbe, bâtiments, vêtements — puis applique des stratégies de traitement différentes à chaque zone.

\begin{itemize}
  \item \textbf{Ciel}: Renforce le bleu et réduit la luminosité séparément pour éviter la surexposition.
  \item \textbf{Visage}: Identifie les zones de couleur de peau, éclaircit et lisse le teint sans toucher aux traits du visage ni aux détails des cheveux.
  \item \textbf{Herbe / feuillage}: Augmente la saturation du vert et améliore la netteté des détails.
  \item \textbf{Flou d'arrière-plan}: En mode portrait, l'IA identifie le contour du sujet, floute l'arrière-plan et garde le sujet net.
\end{itemize}

\subsubsection*{Applications de l'apprentissage profond :}

\begin{itemize}
  \item \textbf{Reconnaissance de scène par IA}: L'appareil photo identifie automatiquement des scènes comme « gastronomie », « animal », « nuit », « document », etc., et optimise les paramètres en conséquence. Par exemple, le mode gastronomie augmente les tons chauds et la saturation pour rendre la nourriture plus appétissante.
  \item \textbf{Flou de portrait par IA}: Les objectifs de smartphone ne peuvent pas physiquement reproduire la faible profondeur de champ d'un reflex ; l'IA utilise un modèle d'apprentissage profond pour prédire les contours du sujet, puis floute l'arrière-plan. Les algorithmes les plus récents peuvent gérer correctement des détails comme les brins de cheveux et les bords en verre.
  \item \textbf{Effacement / extension d'image par IA}: Sélectionnez un passant ou un objet dans l'image, et l'IA « complète » la partie masquée à partir du contexte. C'est une application de l'IA générative dans le domaine de l'imagerie.
\end{itemize}

\subsubsection*{Photographie computationnelle vs photographie optique :}

\begin{itemize}
  \item \textbf{Avantages de la photographie computationnelle}: Bonne qualité d'image même avec un petit objectif, des photos exploitables même en très faible lumière ou en contre-jour, et des photos « belles » sans que l'utilisateur ait besoin de connaissances en photographie.
  \item \textbf{Inconvénients de la photographie computationnelle}: L'image est parfois « surtraitée » — netteté excessive, couleurs déformées, ciel aux teintes peu naturelles ; la fusion multi-images laisse des traînées lorsque le sujet bouge ; les détails sont « calculés » plutôt qu'enregistrés fidèlement.
  \item \textbf{Avantages de la photographie optique}: Un grand capteur avec un objectif à grande ouverture capture la lumière de manière fidèle, qualité d'image naturelle, grand espace de post-traitement ; pas d'artefacts dus à des « extrapolations » algorithmiques.
  \item \textbf{Tendance}: Les smartphones phares modernes combinent « optique + computation » — utilisant un grand capteur pour obtenir des informations réelles et des algorithmes pour optimiser le rendu. Les deux ne sont pas opposés.
\end{itemize}

\booksubsection{VII. Types d'objectifs et focales}

Les smartphones phares modernes sont généralement équipés de trois à quatre objectifs, chacun correspondant à une focale différente. Comprendre les focales vous permet de choisir le bon objectif selon la scène.

\subsubsection*{Objectifs courants et focales équivalentes :}

\begin{itemize}
  \item \textbf{Ultra grand-angle (équivalent 13-24 mm)}: Champ de vision le plus large, idéal pour les paysages, l'architecture et les espaces intérieurs. Capture de vastes scènes, mais les bords présentent une distorsion (les lignes droites deviennent courbes) et les visages sont étirés en portrait. Convient aux « grandes scènes », pas aux « gros plans ».
  \item \textbf{Objectif principal (équivalent 24-28 mm)}: L'objectif le plus utilisé au quotidien, offrant la meilleure qualité d'image et la plus grande ouverture. Équilibre champ de vision et qualité, convenant aux portraits comme aux paysages. Le « 1x » de la plupart des smartphones correspond à l'objectif principal.
  \item \textbf{Téléobjectif (équivalent 70-200 mm)}: Pour les sujets éloignés et les portraits en gros plan. Le téléobjectif comprime la sensation d'espace (les montagnes lointaines et le sujet proche semblent plus rapprochés) et permet des proportions faciales plus naturelles et un meilleur flou d'arrière-plan en portrait.
  \item \textbf{Téléobjectif périscopique (équivalent 200 mm+)}: Grâce à un parcours optique en réflexion, atteint des focales plus longues dans l'épaisseur du smartphone. Permet de capturer des détails à très longue distance, idéal pour les concerts, les événements sportifs et la faune. Réservé aux modèles phares.
\end{itemize}

\subsubsection*{Compromis entre taille de capteur et taille des pixels :}

Le capteur est la « base » qui détermine la quantité de lumière captée. Dans le secteur, on dit qu'« un capteur plus large fait la différence » — à nombre de pixels égal, plus le capteur est grand, plus chaque pixel est grand, plus la lumière captée est importante, et meilleure est la qualité d'image.

\begin{itemize}
  \item \textbf{Grand capteur}: Grande quantité de lumière captée, bonne qualité d'image de nuit, grande plage dynamique. Inconvénients : le smartphone doit être plus épais (module de caméra proéminent) et le coût est élevé.
  \item \textbf{Hauts pixels, petit capteur}: En plein jour avec une lumière abondante, les hauts pixels produisent des photos plus nettes, facilitant le recadrage en post-traitement. Mais en faible lumière, chaque pixel reçoit moins de lumière, avec davantage de bruit.
  \item \textbf{Choix des modèles phares}: L'objectif principal utilise généralement un capteur de 1/1,3 pouce à 1 pouce (certains modèles utilisent déjà un capteur de 1 pouce), avec environ 50 millions de pixels — c'est le point d'équilibre entre qualité d'image et compacité.
\end{itemize}

\subsubsection*{Technologie de fusion de pixels (Pixel Binning) :}

Pour concilier hauts pixels de jour et grands pixels de nuit, les fabricants utilisent la « fusion de pixels » — en combinant plusieurs petits pixels en un seul grand pixel pour la sortie.

\begin{itemize}
  \item \textbf{Pixels quad-bin (Quad Bayer)}: Capteur de 50 millions de pixels, sortie en 12,5 millions de pixels. Chaque pixel de sortie est synthétisé à partir de 4 pixels adjacents, multipliant par 4 la quantité de lumière, pour une qualité d'image plus nette.
  \item \textbf{Pixels nano-bin (Nonacell)}: Utilisé pour les capteurs de plus de 100 millions de pixels, sortie d'environ 11-12 millions de pixels. Même principe, quantité de lumière encore plus grande.
  \item \textbf{Effet réel}: En mode photo par défaut, vous obtenez une photo à basse résolution fusionnée — bonne qualité d'image et fichier réduit. Pour des détails extrêmes (par exemple, documents, sujets éloignés), vous pouvez passer manuellement en « mode hauts pixels » pour une sortie en pleine résolution.
\end{itemize}

\booksubsection{VIII. Modes de prise de vue spéciaux}

Outre la photographie quotidienne, le smartphone propose plusieurs modes spéciaux qui, bien utilisés, permettent de capturer des images impossibles à obtenir autrement.

\subsubsection*{Mode panorama :}

\begin{itemize}
  \item Pendant la prise de vue, déplacez lentement et régulièrement le smartphone dans le sens de la flèche, sans précipitation — un mouvement trop rapide provoque des erreurs d'assemblage.
  \item Gardez le smartphone stable verticalement, déplacez-le uniquement horizontalement. En cas de mouvement vertical, la ligne d'horizon deviendra ondulée.
  \item À utiliser avec prudence en présence de personnes ou de véhicules en mouvement — l'assemblage peut produire des « demi-corps » ou des « fantômes ».
  \item Le mode panorama peut aussi être utilisé en vertical — prise de vue verticale et assemblage horizontal, idéal pour les bâtiments hauts ou les cascades.
\end{itemize}

\subsubsection*{Ralenti et accéléré :}

\begin{itemize}
  \item \textbf{Ralenti}: Filmez à 120 fps ou 240 fps en haute fréquence d'images, la lecture devient un ralenti. Idéal pour les gouttes d'eau, les instants sportifs et les sauts d'animaux. La lumière doit être suffisante ; en mode ralenti, le temps d'exposition est court, la qualité en faible lumière est médiocre.
  \item \textbf{Accéléré (Time-lapse)}: Prend des photos à intervalles réguliers et les assemble en une vidéo accélérée. Idéal pour les flux de voitures, le mouvement des nuages, les couchers de soleil et l'éclosion des fleurs. La durée de prise de vue varie de quelques minutes à plusieurs heures ; utilisez impérativement un trépied, à main levée ce sera flou.
  \item \textbf{Astuce}: Pour filmer les nuages en accéléré, fixez le smartphone et filmez pendant 15-30 minutes ; vous obtiendrez une vidéo accélérée d'environ 10 secondes, avec un effet saisissant.
\end{itemize}

\subsubsection*{Mode étoiles / astronomie :}

\begin{itemize}
  \item Il faut s'éloigner de la pollution lumineuse urbaine et trouver une zone rurale ou montagneuse sans éclairage public. La Voie lactée ne peut pas être photographiée en ville.
  \item L'utilisation d'un trépied est obligatoire, à main levée ce sera flou. Le smartphone doit être maintenu immobile pendant plusieurs secondes à plusieurs dizaines de secondes.
  \item Le temps doit être dégagé sans nuages, la Lune ne doit pas être trop grosse (la pleine Lune masque les étoiles). Les nuits sombres du début ou de la fin du mois lunaire sont les plus appropriées.
  \item Le « Mode nuit » de l'iPhone reconnaît automatiquement le ciel étoilé et prolonge le temps d'exposition (jusqu'à 30 secondes) ; les « Mode étoiles » ou « Mode astronomie » des smartphones Android phares offrent des résultats similaires. Soyez patient et ne bougez pas le smartphone après avoir appuyé sur le déclencheur.
\end{itemize}

\subsubsection*{Mode macro :}

\begin{itemize}
  \item Pour photographier de petits objets tels que fleurs, insectes, bijoux et gros plans sur du texte. Mise au point possible à 1-2 cm de distance.
  \item Certains smartphones disposent d'un objectif macro dédié (souvent un petit capteur de 2 ou 5 mégapixels), avec une qualité moyenne ; les modèles phares utilisent l'ultra grand-angle pour la macro, avec une meilleure qualité.
  \item En prise de vue rapprochée, l'objectif est très près du sujet et peut bloquer la lumière ; utilisez une lampe de bureau ou une lampe de poche pour un éclairage d'appoint si nécessaire.
  \item La distance de mise au point est très courte, la profondeur de champ est très faible ; un déplacement du sujet de quelques millimètres peut provoquer une mise au point manquée. Tenez fermement et de préférence immobilisez le sujet.
\end{itemize}

\begin{tcbinfo}{Conseils pratiques pour la photographie statique}
\begin{itemize}
  \item Photographie quotidienne : HEIF/JPEG par défaut suffit, fichiers petits et qualité adéquate.
  \item Scènes importantes (paysages, portraits, post-traitement nécessaire) : activez le RAW pour conserver un maximum d'espace de post-traitement.
  \item Choisissez le bon objectif : ultra grand-angle pour les paysages, téléobjectif pour les portraits, objectif principal pour l'usage quotidien.
  \item Nuit et étoiles : fixez impérativement le smartphone, activez le mode nuit ou le mode astronomie.
  \item Les smartphones à grand capteur (par exemple avec un capteur de 1 pouce) offrent un plafond de qualité plus élevé, mais les habitudes de prise de vue restent déterminantes.
\end{itemize}

\end{tcbinfo}

\begin{tcbpractice}{À essayer : Explorez les capacités de votre appareil photo}
\begin{enumerate}
  \item Vérifiez si votre smartphone prend en charge la prise de vue RAW (Réglages Appareil photo Formats), activez-la et prenez une photo, puis comparez la taille des fichiers JPEG et RAW de la même scène.
  \item Passez aux trois focales ultra grand-angle, objectif principal et téléobjectif, photographiez le même paysage éloigné et comparez les différences de champ de vision et de qualité d'image.
  \item Trouvez une scène avec un mode nuit ou un mode étoiles, fixez le smartphone et prenez une photo de nuit.
  \item Essayez de prendre un panorama en mode horizontal, en veignant à déplacer le smartphone de manière fluide.
\end{enumerate}

\end{tcbpractice}
